

\documentclass{article}

% -----------------------------
% Packages
% -----------------------------
\usepackage{fancyhdr}
\usepackage{extramarks}
\usepackage{amsmath}
\usepackage{amsthm}
\usepackage{amsfonts}
\usepackage{tikz}
\usepackage[plain]{algorithm}
\usepackage{algpseudocode}
\usepackage{amssymb}

\usetikzlibrary{automata,positioning}

% -----------------------------
% Basic Document Settings
% -----------------------------
\topmargin=-0.45in
\evensidemargin=0in
\oddsidemargin=0in
\textwidth=6.5in
\textheight=9.0in
\headsep=0.25in

\linespread{1.1}

\pagestyle{fancy}
\lhead{\hmwkAuthorName}
\chead{\hmwkClass: \hmwkTitle}
\rhead{}
\lfoot{\lastxmark}
\cfoot{\thepage}

\renewcommand\headrulewidth{0.4pt}
\renewcommand\footrulewidth{0.4pt}

\setlength\parindent{0pt}

% -----------------------------
% Problem Section Helpers
% -----------------------------
\newcommand{\enterProblemHeader}[1]{
  \nobreak\extramarks{}{Problem #1 continued on next page\ldots}\nobreak{}
  \nobreak\extramarks{Problem #1 (continued)}{Problem #1 continued on next page\ldots}\nobreak{}
}

\newcommand{\exitProblemHeader}[1]{
  \nobreak\extramarks{Problem #1 (continued)}{Problem #1 continued on next page\ldots}\nobreak{}
  \nobreak\extramarks{Problem #1}{}\nobreak{}
}

\setcounter{secnumdepth}{0}
\newcounter{partCounter}
\newcounter{homeworkProblemCounter}
\setcounter{homeworkProblemCounter}{1}
\nobreak\extramarks{Problem \arabic{homeworkProblemCounter}}{}\nobreak{}

% -----------------------------
% Theorem env (optional)
% -----------------------------
\newtheorem*{theorem}{Theorem}

% -----------------------------
% Homework Problem Environment
%   Usage:
%   \begin{homeworkProblem}            % auto-number
%     % Your answer here
%   \end{homeworkProblem}
%
%   \begin{homeworkProblem}[1.2.3]     % label problem "1.2.3" (e.g. textbook numbering)
%     % Your answer here
%   \end{homeworkProblem}
% -----------------------------
\newenvironment{homeworkProblem}[1][]{
  \def\hwProblemLabel{#1}
  \ifx\hwProblemLabel\empty
    \edef\hwProblemLabel{\arabic{homeworkProblemCounter}}
    \stepcounter{homeworkProblemCounter}
  \fi
  \section{Problem \hwProblemLabel}
  \setcounter{partCounter}{1}
  \enterProblemHeader{\hwProblemLabel}
}{
  \exitProblemHeader{\hwProblemLabel}
}

% -----------------------------
% Homework Details (EDIT THESE)
% -----------------------------
\newcommand{\hmwkTitle}{Homework \#4}
\newcommand{\hmwkClass}{Math 3320}
\newcommand{\hmwkAuthorName}{\textbf{Joseph Quinn}}

% -----------------------------
% Title
% -----------------------------
\title{
  \vspace{2in}
  \textmd{\textbf{\hmwkClass:\ \hmwkTitle}}\\
  \vspace{3in}
}
\author{\hmwkAuthorName}
\date{}

% -----------------------------
% Convenience Macros (optional)
% -----------------------------
\renewcommand{\part}[1]{\textbf{\large Part \Alph{partCounter}}\stepcounter{partCounter}\\}

% Algorithms
\newcommand{\alg}[1]{\textsc{\bfseries \footnotesize #1}}

% Calculus
\newcommand{\deriv}[1]{\frac{\mathrm{d}}{\mathrm{d}x} (#1)}
\newcommand{\pderiv}[2]{\frac{\partial}{\partial #1} (#2)}
\newcommand{\dx}{\mathrm{d}x}

% Statistics
\newcommand{\E}{\mathrm{E}}
\newcommand{\Var}{\mathrm{Var}}
\newcommand{\Cov}{\mathrm{Cov}}
\newcommand{\Bias}{\mathrm{Bias}}

% Proof step (for aligned derivations)
\newcommand{\step}[2]{& #1 & & \text{#2} \\}

% Solution header (optional)
\newcommand{\solution}{\textbf{\large Solution}}

% -----------------------------
% Document
% -----------------------------
\begin{document}

\maketitle
\pagebreak

\begin{homeworkProblem}[A5]
	\begin{align*}
		\text{max} \quad  & = 13x_{1} + 3x_{2} + 12x_{3} + 108x_{4} \\
		\text{s.t.} \quad & = x_{1}+2x_{2}+3x_{3}+27x_{4} \le 24    \\
		                  & = 6x_{1}+x_{2}+4x_{3}+36x_{4} \le 28    \\
		                  & = x_{1},x_{2},x_{3},x_{4} \ge 0
	\end{align*}

	\begin{enumerate}
		\item[(I1)]
		      This is already in standard form so we can proceed
		\item[(I2)]
		      Setup the initial dictionary like so,
		      \begin{align*}
			      x_{5} = & 24 - x_{1}-2x_{2}-3x_{3}-27x_{4} \\
			      x_{6} = & 28-6x_{1}-x_{2}-4x_{3}-36x_{4}   \\
			      z =     & 13x_{1}+3x_{2}+12x_{3}+108x_{4}
		      \end{align*}
		\item[(I3)]
		      For this dictionary the basic feasible solution where $(x_{1} = 0, x_{2}=0, x_{3}=0, x_{4}=0, x_{5}=24, x_{6}=28)$
		      with $z=0$

		\item[(P1)]
		      According to Bland's rule the entering variable from the positive coefficients, is $x_{1}$
		\item[(P2)]
		      Via the ratio test we see,
		      \begin{align*}
			      x_{1} \le 24  \text{ for } x_{5}          \\
			      x_{1} \le \frac{28}{6} \text{ for } x_{6} \\
		      \end{align*}
		      So $x_{6}$ has the tightest bounds and is our leaving variable.

		\item[(P3)]
		      After performing the pivot, we get the dictionary,

		      \begin{align*}
			      x_{5} = & \frac{58}{3} - \frac{11}{6}x_{2} - \frac{7}{3}x_{3} - 21x_{4} + \frac{1}{6}x_{6}   \\
			      x_{1} = & \frac{14}{3} - \frac{1}{6}x_{2} - \frac{2}{3}x_{3} - 6x_{4} - \frac{1}{6}x_{6}     \\
			      z =     & \frac{182}{3} + \frac{5}{6}x_{2} + \frac{10}{3}x_{3} + 30x_{4} - \frac{13}{6}x_{6}
		      \end{align*}

		\item[(P4)]
		      For this dictionary the basic feasible solution where $(x_{1} = \frac{14}{3}, x_{2}=0, x_{3}=0, x_{4}=0, x_{5}=\frac{58}{3}, x_{6}=0)$
		      with $z=\frac{182}{3}$
		\item[(P1)]
		      According to Bland's rule the entering variable from the positive coefficients, is $x_{2}$
		\item[(P2)]
		      \begin{align*}
			      x_{2} \le \frac{116}{11} \text{ for } x_{5} \\
			      x_{2} \le 28 \text{ for } x_{1}             \\
		      \end{align*}

		      This makes the tightest constraints $x_{5}$ and it becomes our leaving variable
		\item[(P3)]
		      Performing the pivot results in the dictionary,
		      \begin{align*}
			      x_{2} = & \frac{116}{11} - \frac{14}{11}x_{3} - \frac{126}{11}x_{4} - \frac{6}{11}x_{5} + \frac{1}{11}x_{6}  \\
			      x_{1} = & \frac{32}{11} - \frac{5}{11}x_{3} - \frac{45}{11}x_{4} + \frac{1}{11}x_{5} - \frac{2}{11}x_{6}     \\
			      z =     & \frac{764}{11} + \frac{25}{11}x_{3} + \frac{225}{11}x_{4} - \frac{5}{11}x_{5} - \frac{23}{11}x_{6}
		      \end{align*}

		\item[(P4)]
		      This results in the basic feasible solution with $(x_{1}=\frac{14}{3},x_{2}=0, x_{4}=0,x_{4}=0,x_{5}=\frac{58}{3},x_{6}=0)$
		      resulting in $z=\frac{182}{3}$

		\item[(P1)]
		      Based on blands rule, our entering vairable is $x_{3}$
		\item[(P2)]
		      Via the ratio test
		      \begin{align*}
			      x_{3} \le \frac{116}{14} \text{ for } x_{2} \\
			      x_{3} \le \frac{32}{5} \text{ for } x_{1}   \\
		      \end{align*}

		\item[(P3)]
		      After pivoting this results in the dictionary,
		      \begin{align*}
			      x_{2} = & \frac{12}{5} + \frac{14}{5}x_{1} - \frac{4}{5}x_{5} + \frac{3}{5}x_{6}          \\
			      x_{3} = & \frac{32}{5} - \frac{11}{5}x_{1} - 9x_{4} + \frac{1}{5}x_{5} - \frac{2}{5}x_{6} \\
			      z =     & 84 - 5x_{1} - 3x_{6}
		      \end{align*}

		      Now we notice that we have no more coefficents with positive values in the $z$ equation meaning we cannot increase $z$ any further
		      aka we have found our optimal solution.

		      However we see that we have some nonbasic vairables that have 0 coefficents, meaning there values can vary resulting in the same $z$ value.

		      and since none of our contraints have a 0 value contant we can vary these values and remain a feasible solution.

		      This means our optimal solution is the set of solutions,
		      \[
			      x_1 = 0,\qquad
			      x_2 = \frac{12}{5} - \frac{4}{5}x_5,\qquad
			      x_3 = \frac{32}{5} - 9x_4 + \frac{1}{5}x_5,
		      \]

		      where \(x_4\) and \(x_5\) vary, subject to

		      \[
			      x_4 \ge 0,\qquad
			      x_5 \ge 0,\qquad
			      x_5 \le 3,\qquad
			      45x_4 - x_5 \le 32.
		      \]

	\end{enumerate}
\end{homeworkProblem}

\end{document}

